\section{引言}
\label{sec:intro}

科学出版物的快速增长给研究人员开展系统性文献综述带来了巨大挑战。据估计，每年各学科发表的新学术论文超过500万篇 \citep{bornmann-etal-2021}。一篇系统性文献综述通常要求研究人员手动筛选数千篇论文、识别相关主题并综合研究结果——这一过程可能耗时数周甚至数月。这种劳动密集型过程不仅耗时，还容易受到选择偏见和覆盖不全的影响。

近年来，大语言模型（LLM）和自然语言处理技术的进展为自动化文献综述的各个阶段开辟了新的可能。然而，现有方法往往只产生浅层摘要，仅仅罗列方法名称，而不解释底层原理、数学公式或性能差异背后的物理原因。一个真正有用的自动化综述应该超越表面描述，提供公式级分析、推导链条和第一性原理推理。

本文提出了一套多Agent自动化文献综述系统，通过以下四项关键创新解决上述挑战：

\begin{enumerate}
    \item \textbf{协调者-代理架构}：由四个专业化Agent（检索、筛选、聚类、摘要）组成顺序流水线，以层级研究问题（RQ）树驱动，通过检查点恢复机制保证数据流的连贯性和质量控制。
    \item \textbf{三级筛选机制}：结合基于嵌入的余弦相似度与LLM辅助的边界区域判定，在相关性的吞吐量和精度之间取得平衡。
    \item \textbf{GPU加速的语义聚类}：采用PCA降维和HDBSCAN聚类，支持自适应参数调节和领域一致性过滤，生成有意义的主题簇。
    \item \textbf{聚焦原理的报告生成}：通过反幻觉提示词工程强制要求公式解读、推导链条、理论极限分析和严格引用验证，生成分析"方法为何有效"而非仅仅罗列"方法是什么"的综述报告。
\end{enumerate}

系统集成五个免费学术数据库——arXiv、PubMed、DBLP、Europe PMC和OpenAlex——基本操作无需API密钥。支持命令行和Web界面两种交互方式，便于各学科研究人员使用。

本文其余部分组织如下：第~\ref{sec:lit-rev}~节综述相关工作；第~\ref{sec:model}~节描述系统架构和形式化问题定义；第~\ref{sec:method}~节详述各Agent的技术实现；第~\ref{sec:exp}~节展示实验结果和分析；第~\ref{sec:concl}~节总结全文。
